\subsection{Module 8: RasC--TORC2--PKB--CRAC Relay}\label{sec:m8}

\paragraph{Role.}
The relay of the cAMP signal, i.e.\ the re-synthesis of cAMP in response to cAMP, runs through a second Ras branch.
RasC and RasG are both activated by cAMP, but RasC is responsible for adenylyl cyclase activation and RasG for chemotaxis; RasGEFA activates RasC but not RasG~\cite{kae2007}.
RasC is an upstream regulator of the target of rapamycin complex 2 (TORC2), which phosphorylates the myristoylated kinase PKBR1 independently of PIP3 and prolongs the time course of ACA activation~\cite{cai2010}.
Within seconds of a stimulus, TORC2 activates two PKB homologues, PkbA and PKBR1, also in cells lacking PI3K activity~\cite{kamimura2008}.
PKBR1 has no PH domain, is myristoylated for constitutive membrane localisation and does not require PI3K~\cite{meili2000}.
PKB and PKBR1 phosphorylate the scaffold Sca1 of a Sca1/RasGEF/PP2A complex and thereby control RasC in a negative feedback~\cite{charest2010}, and PKA exerts a further negative feedback on Ras, Rap1 and TORC2 signalling~\cite{scavello2017}.
The effector CRAC (a PH-domain protein) is required for receptor- and G-protein-mediated activation of ACA~\cite{insall1994}; CRAC binds PI3K products~\cite{comer2005} and is the translocation reporter of the G-protein pathway~\cite{parent1998}.
The model connects this chain to the excitable element of Module~0 by an output edge (8.14) that \emph{injects the excitable activator X}, and not by a direct activation of ACA.
A direct receptor-to-ACA edge would make the receptor arm a second pulse generator with its own amplitude and refractory state, whereas injecting X makes a relayed pulse identical to a spontaneous one and lets one refractory clock (Module~0) gate both.

\begin{table*}[!t]
\centering
\caption{Module~8 reactions: the RasC$\to$TORC2$\to$PKB$\to$CRAC relay arm and its output edge 8.14 into the excitable element of Module~0. PKBR1 (module 8b) and CRAC are the kinase and effector; $g(\mathrm{PIP3})=[1+(K_P/N_\mathrm{PIP3})^{8}]^{-1}$ and $b(\mathrm{Rf})$ are the gates of Sec.~\ref{sec:hill}.}
\label{tab:m8rxn}
\footnotesize
\renewcommand{\arraystretch}{1.15}
\begin{tabularx}{\textwidth}{@{}l>{\raggedright\arraybackslash}p{0.37\textwidth}>{\raggedright\arraybackslash}X@{}}
\toprule
ID & Reaction & Propensity (rate law)\\
\midrule
8.1 & $\mathrm{G}_{\beta\gamma}+\mathrm{GEF_A} \to \mathrm{GEF_A^*}+\mathrm{G}_{\beta\gamma}$ & $k_{A,\mathrm{on}}[\mathrm{G}_{\beta\gamma}][\mathrm{GEF_A}]$ \\
8.2 & $\mathrm{GEF_A^*} \to \mathrm{GEF_A}$ & $k_{A,\mathrm{off}}[\mathrm{GEF_A^*}]$ \\
8.2s & $\mathrm{B}^*+\mathrm{GEF_A^*} \to \mathrm{GEF_A}+\mathrm{B}^*$ & $k_S[\mathrm{B}^*][\mathrm{GEF_A^*}]$ \\
8.3 & $\mathrm{GEF_A^*}+\mathrm{RasC}^{\mathrm{GDP}} \to \mathrm{RasC}^{\mathrm{GTP}}+\mathrm{GEF_A^*}$ & $k_{C,\mathrm{on}}[\mathrm{GEF_A^*}][\mathrm{RasC}^{\mathrm{GDP}}]$ \\
8.4 & $\mathrm{RasC}^{\mathrm{GTP}} \to \mathrm{RasC}^{\mathrm{GDP}}$ & $k_{C,\mathrm{off}}[\mathrm{RasC}^{\mathrm{GTP}}]$ \\
8.4b & $\mathrm{RC}+\mathrm{CGAP} \to \mathrm{RC}+\mathrm{CGAP^*}$ & $k_{CG,\mathrm{on}}[\mathrm{RC}][\mathrm{CGAP}]$ \\
8.4c & $\mathrm{CGAP^*} \to \mathrm{CGAP}$ & $k_{CG,\mathrm{off}}[\mathrm{CGAP^*}]$ \\
8.4d & $\mathrm{CGAP^*}+\mathrm{RasC}^{\mathrm{GTP}} \to \mathrm{RasC}^{\mathrm{GDP}}+\mathrm{CGAP^*}$ & $k_{CG,\mathrm{hyd}}[\mathrm{CGAP^*}][\mathrm{RasC}^{\mathrm{GTP}}]$ \\
8.5 & $\mathrm{RasC}^{\mathrm{GTP}}+\mathrm{TORC2} \to \mathrm{TORC2^*}+\mathrm{RasC}^{\mathrm{GTP}}$ & $k_{T,\mathrm{on}}[\mathrm{RasC}^{\mathrm{GTP}}][\mathrm{TORC2}]$ \\
8.6 & $\mathrm{TORC2^*} \to \mathrm{TORC2}$ & $k_{T,\mathrm{off}}[\mathrm{TORC2^*}]$ \\
8b.1 & $\mathrm{TORC2^*}+\mathrm{PKBR1} \to \mathrm{PKBR1^*}+\mathrm{TORC2^*}$ & $k_{P,\mathrm{on}}[\mathrm{TORC2^*}][\mathrm{PKBR1}]$ \\
8b.2 & $\mathrm{PKBR1^*} \to \mathrm{PKBR1}$ & $k_{P,\mathrm{off}}[\mathrm{PKBR1^*}]$ \\
8b.3 & $\mathrm{PKBR1^*}+\mathrm{CRAC}_m \to \mathrm{CRAC}_m^*+\mathrm{PKBR1^*}$ & $k_\mathrm{Cr}[\mathrm{PKBR1^*}][\mathrm{CRAC}_m]$ \\
8b.4 & $\mathrm{PKA^*}+\mathrm{PKBR1^*} \to \mathrm{PKA^*}+\mathrm{PKBR1}$ & $k_{PKA}[\mathrm{PKA^*}][\mathrm{PKBR1^*}]$ \\
8.9 & $\mathrm{PKBA}_m^*+\mathrm{CRAC}_m \to \mathrm{CRAC}_m^*+\mathrm{PKBA}_m^*$ & $k_\mathrm{Cr}[\mathrm{PKBA}_m^*][\mathrm{CRAC}_m]$ \\
8.9b & $\mathrm{PIP3}+\mathrm{CRAC}_c \to \mathrm{CRAC}_m+\mathrm{PIP3}$ & $k_{\mathrm{Cr},P}[\mathrm{PIP3}][\mathrm{CRAC}_c]$ \\
8.9c & $\mathrm{CRAC}_c \to \mathrm{CRAC}_m$ & $k_{\mathrm{Cr},b}[\mathrm{CRAC}_c]$ \\
8.10 & $\mathrm{CRAC}_m^* \to \mathrm{CRAC}_m$ & $k_{\mathrm{Cr},\mathrm{dp}}[\mathrm{CRAC}_m^*]$ \\
8.10b & $\mathrm{CRAC}_m \to \mathrm{CRAC}_c$ & $k_{\mathrm{Cr},\mathrm{off}}[\mathrm{CRAC}_m]$ \\
8.14 & $\mathrm{CRAC}_m^* \to \mathrm{CRAC}_m^*+\mathrm{X}$ & $\dfrac{k_{X}}{N_\mathrm{CRAC}}\,N_{\mathrm{CRAC}^*}\,g(\mathrm{PIP3})\,b(\mathrm{Rf})\quad[\text{molec\,s}^{-1}\text{cell}^{-1}]$ \\
\bottomrule
\end{tabularx}
\end{table*}

\begin{table*}[!t]
\centering
\caption{Module~8 constants (a): RasC branch, TORC2 and PKBR1.}
\label{tab:m8const}
\footnotesize
\renewcommand{\arraystretch}{1.15}
\begin{tabularx}{\textwidth}{@{}l r l c >{\raggedright\arraybackslash}X@{}}
\toprule
Symbol & Value & Unit & Basis & Meaning, justification, source\\
\midrule
$[\mathrm{GEF_A}]_\mathrm{tot}$ & \num{1600} & $\mathrm{molec/voxel}$ & D & Sized so that the number of active RasGEFA at the driven state (67\,000 G$\beta\gamma$/voxel) equals that of the previous, saturated parameterisation. \\
$k_{A,\mathrm{on}}$ & \num{7.6e-4} & $\mu\mathrm{M}^{-1}\mathrm{s}^{-1}$ & D & $k_{E,\mathrm{on}}k_{A,\mathrm{off}}/k_{E,\mathrm{off}}$: RasGEFA and RasGEFR are both G$\beta\gamma$-driven GEFs~\cite{kae2007} and share the activated-fraction curve in G$\beta\gamma$. \\
$k_{A,\mathrm{off}}$ & \num{0.1} & $\mathrm{s^{-1}}$ & E & RasGEFA inactivation. \\
$k_S$ & \num{1} & $\mu\mathrm{M}^{-1}\mathrm{s}^{-1}$ & C & Release of active RasGEFA by the PKB-phosphorylated Sca1-like brake (negative feedback on RasC); calibrated so that RasC returns near baseline within $\approx\!40$\,s~\cite{charest2010}. \\
$k_{C,\mathrm{on}}$ & \num{2} & $\mu\mathrm{M}^{-1}\mathrm{s}^{-1}$ & E & RasGEFA-catalysed exchange on RasC~\cite{kae2007}. \\
$k_{C,\mathrm{off}}$ & \num{0.02} & $\mathrm{s^{-1}}$ & E & Intrinsic RasC hydrolysis (constitutive, slow). \\
$k_{CG,\mathrm{on}},\ k_{CG,\mathrm{off}},\ k_{CG,\mathrm{hyd}}$ & \num{0.05},\ \num{0.005},\ \num{150} & $\mu\mathrm{M}^{-1}\mathrm{s}^{-1},\ \mathrm{s^{-1}},\ \mu\mathrm{M}^{-1}\mathrm{s}^{-1}$ & E & Receptor-driven, slowly decaying RasC GAP that holds RasC adapted. \\
$k_{T,\mathrm{on}}$ & \num{50} & $\mu\mathrm{M}^{-1}\mathrm{s}^{-1}$ & C & TORC2 activation by RasC-GTP~\cite{cai2010}. \\
$k_{T,\mathrm{off}}$ & \num{0.5} & $\mathrm{s^{-1}}$ & C & TORC2 inactivation ($\tau=2$\,s). $k_{T,\mathrm{on}}$ and $k_{T,\mathrm{off}}$ were scaled together, which leaves the activated fraction unchanged and shortens only the latency, to a peak at 5--10\,s as observed for phosphorylation of Sca1~\cite{charest2010}. Other PKB substrates and the PKBR1 hydrophobic motif peak later, at 20--60\,s~\cite{kamimura2008,cai2010}, so the model is on the fast side. \\
$k_{P,\mathrm{on}},\ k_{P,\mathrm{off}}$ & \num{30},\ \num{0.5} & $\mu\mathrm{M}^{-1}\mathrm{s}^{-1},\ \mathrm{s^{-1}}$ & E & PKBR1 activation by TORC2 and its dephosphorylation ($\tau=2$\,s); PKBR1 is constitutively membrane-bound and PI3K-independent~\cite{meili2000,kamimura2008}. \emph{Departure:} the measured half-life of PKBR1 hydrophobic-motif phosphorylation is about 40\,s~\cite{cai2010}. \\
\bottomrule
\end{tabularx}
\end{table*}

\begin{table*}[!t]
\centering
\caption{Module~8 constants (b): CRAC, the relay output edge, pools and diffusion constants.}
\label{tab:m8constb}
\footnotesize
\renewcommand{\arraystretch}{1.15}
\begin{tabularx}{\textwidth}{@{}l r l c >{\raggedright\arraybackslash}X@{}}
\toprule
Symbol & Value & Unit & Basis & Meaning, justification, source\\
\midrule
$k_\mathrm{Cr}$ & \num{10} & $\mu\mathrm{M}^{-1}\mathrm{s}^{-1}$ & C & Phosphorylation of membrane-bound CRAC by either PKB. \mkt{M8-tor-pkbs}{Both PKBs are TORC2 targets}~\cite{kamimura2008} and \mkt{M8-aca-prolong}{RasC--TORC2--PKBR1 signalling prolongs ACA activation}~\cite{cai2010}; PKB$\to$CRAC is a model assumption. \\
$k_{\mathrm{Cr},P}$ & \num{1} & $\mu\mathrm{M}^{-1}\mathrm{s}^{-1}$ & E & PIP3-facilitated recruitment of CRAC, which is \mkt{M8-crac-pip3}{a PH-domain protein}~\cite{insall1994,parent1998,comer2005}. \\
$k_{\mathrm{Cr},b}$ & \num{0.2} & $\mathrm{s^{-1}}$ & C & PIP3-independent membrane association; keeps the relay functional when PIP3 is absent (\mkt{M8-pi3k-ind}{PI3K-independent PKB signalling}~\cite{kamimura2008}). \\
$k_{\mathrm{Cr},\mathrm{off}}$ & \num{0.1} & $\mathrm{s^{-1}}$ & E & Release of unphosphorylated CRAC from the membrane. \\
$k_{\mathrm{Cr},\mathrm{dp}}$ & \num{1.5} & $\mathrm{s^{-1}}$ & D & $=k_\mathrm{Cr}\,a\,(1-f)/f$ with $a=567$ PKBA$^*$/voxel (the measured peak) and target CRAC$^*$ fraction $f=0.44$; fixes the activated fraction of the CRAC switch. \\
$k_X$ & \num{200} & $\mathrm{molec\,s^{-1}\,cell^{-1}}$ & C & Injection rate of X at full CRAC$^*$ (per-CRAC$^*$ rate $k_X/N_\mathrm{CRAC}=200/2002=0.1\,\mathrm{s^{-1}}$, with $N_\mathrm{CRAC}=2002$ per cell); must exceed the $\approx\!24$-molecule threshold of Module~0. \\
$K_P,\ n_P$ & \num{5000},\ 8 & $\mathrm{molec/voxel},\ -$ & C & Half-open level and steepness of the PIP3 gate (Sec.~\ref{sec:hill}). \\
$k_\mathrm{PKA}$ & \num{0.938} & $\mu\mathrm{M}^{-1}\mathrm{s}^{-1}$ & E & PKA$^*$ feedback on PKBR1$^*$; \mkt{M8-pka-fb}{PKA exerts negative feedback on Ras, Rap1 and TORC2 signalling}~\cite{scavello2017}. The magnitude is a design choice: the value $0.6\,\mathrm{min^{-1}}$ is \mkt{M8-laub-k}{the coefficient $k_5$ of}~\cite{laub1998} (CAR1$\to$ERK2, not a PKA coefficient; the PKA$\to$ERK2 coefficient is $k_6=0.8\,\mathrm{min^{-1}}$, used in Module~10), $30\times$ boosted and normalised by the PKA pool. \\
$[\mathrm{RasC}]_\mathrm{tot}$ & \num{385} & $\mathrm{molec/voxel}$ & E & Pool size; not constrained by measurements. \\
$[\mathrm{CGAP}]_\mathrm{tot},\ [\mathrm{TORC2}]_\mathrm{tot},\ [\mathrm{CRAC}]_\mathrm{tot}$ & \num{154} & $\mathrm{molec/voxel}$ & E & Pool sizes (2\,002 per cell); not constrained by measurements. \\
$[\mathrm{PKBR1}]_\mathrm{tot}$ & \num{77} & $\mathrm{molec/voxel}$ & E & PKBR1 pool; PkbA (Module~4) is a separate, larger pool. \\
$D_{\mathrm{RasC}},D_{\mathrm{TORC2}},D_{\mathrm{PKBR1}},D_{\mathrm{CRAC}_m}$ & \num{0.1} & $\mu\mathrm{m^2/s}$ & E & Membrane-bound species. \\
$D_{\mathrm{GEF_A}},D_{\mathrm{CGAP}},D_{\mathrm{CRAC}_c}$ & \num{10} & $\mu\mathrm{m^2/s}$ & E & Cytosolic value ($D_{\mathrm{GEF_A^*}}=2$). \\
\bottomrule
\end{tabularx}
\end{table*}


\paragraph{Reactions.}
\begin{description}\itemsep0pt
\item[8.1/8.2] RasGEFA is activated by G$\beta\gamma$ and inactivated with $\tau=10$\,s. It uses the same G$\beta\gamma$-dependent activated-fraction curve as the RasG excitor 3.1, since both are G$\beta\gamma$-driven RasGEFs~\cite{kae2007}.
\item[8.2s] The PKB-phosphorylated Sca1-like brake substrate B$^*$ (Module~3) releases active RasGEFA, which is the negative feedback of~\cite{charest2010}. Without PKB or TORC2, RasC activation is increased and fails to adapt by 40\,s; the constant $k_S$ is calibrated to return RasC near baseline within about 40\,s.
\item[8.3/8.4] Active RasGEFA converts RasC-GDP to RasC-GTP; RasC-GTP hydrolyses slowly.
\item[8.4b--d] A receptor-driven, slowly decaying RasC GAP (RC $\to$ CGAP$^*$, decay $\tau=200$\,s) hydrolyses RasC-GTP and holds RasC in its adapted state.
\item[8.5/8.6] RasC-GTP activates TORC2~\cite{cai2010}, which relaxes with $\tau=2$\,s. The short lifetime is a design choice that places the peak of TORC2 targets within 5--10\,s, as observed for the PKB substrate Sca1~\cite{charest2010}.
\item[8b.1/8b.2] TORC2$^*$ phosphorylates the membrane-anchored PKBR1 directly (no recruitment step), and PKBR1$^*$ is dephosphorylated rapidly, so that PKBR1 is a fast-on/fast-off pulse generator that is cleared before the slow negative feedbacks accumulate. Measured PKBR1 hydrophobic-motif phosphorylation is slower: it peaks at 30--60\,s, has a half-life of about 40\,s and is back at baseline after 2--3\,min~\cite{cai2010}; the model's $\tau=2$\,s is a design choice for the fast relay, not a measured value.
\item[8b.4] PKA$^*$ inactivates PKBR1$^*$, which closes a slow negative feedback from cAMP to the relay arm~\cite{scavello2017}.
\item[8.9, 8b.3] Both PKBs phosphorylate the membrane-bound CRAC pool, PkbA (Module~4) through 8.9 and PKBR1 through 8b.3.
\item[8.9b/8.9c/8.10/8.10b] CRAC is a three-state moiety: cytosolic CRAC$_c$, membrane-bound CRAC$_m$ and phosphorylated CRAC$_m^*$. Two facts must hold together and are compatible only if recruitment and activation are separate steps: CRAC is a PH-domain protein recruited by PIP3~\cite{parent1998,comer2005}, and PKB signalling and the relay do not require PIP3~\cite{kamimura2008}. PIP3 therefore recruits CRAC (8.9b), a PIP3-independent association keeps the relay functional without PIP3 (8.9c), and the kinases phosphorylate CRAC on the membrane (8.9, 8b.3). Dephosphorylation (8.10) returns CRAC$^*$ to the membrane pool and not to the cytosol, because a phosphatase does not strip the lipid anchor; releasing it to the cytosol would cap the activated fraction at $k_\mathrm{rec}/(k_\mathrm{rec}+k_\mathrm{off})$, whatever the kinase drive. Release from the membrane (8.10b) is a separate step.
\item[8.14] CRAC$^*$ injects X into the excitable element at up to $k_X=200$ molecules\,s$^{-1}$ per cell, gated by PIP3 and by the refractory clock. The PIP3 gate requires that the cell has already resolved the direction of the signal before it re-amplifies it: PIP3 is the first species of the cascade that is both stimulus-driven and polarised.
\end{description}
Three further edges exist in the code and are excluded from the default network:
the activation of PkbA by TORC2 without a recruitment step (8.7/8.8) is PKBR1's mechanism applied to PkbA's species and would double count the kinase that Module~8b already provides;
the PKA feedback on PkbA (8.13) is excluded together with them;
and the direct activation of ACA by CRAC$^*$ (8.11) is disabled ($k=0$) for the reason stated above.

\paragraph{Constants.}
Most pool sizes and the RasC-branch rate constants are estimates for which no measurements exist; the calibrated properties are the timing (TORC2$^*$ and PKBR1$^*$ peak within 5--10\,s), the activated fraction of the CRAC switch and the relay drive.
Two constants are derived from measurements of other modules and must be re-derived when those change.
First, $k_{\mathrm{Cr},\mathrm{dp}}$ sets the activated fraction of the CRAC phospho-cycle, which is a function of $a/k_{\mathrm{Cr},\mathrm{dp}}$ only, where $a=k_\mathrm{Cr}[\mathrm{PKBA}^*]$; it is derived from the measured peak of PkbA$^*$ (567 molecules per voxel) and the target CRAC$^*$ fraction of 44\,\%.
When Module~4 was re-anchored the measured peak fell tenfold, the fraction fell to 11\,\% and the relay stopped firing, so the derivation is stored with the measurement it depends on.
Second, the half-open level $K_P$ of the PIP3 gate is placed between the PIP3 peak evoked by 2\,nM (about 5000 molecules per voxel) and the standing PIP3 (about $10^3$).
Neither the gain $k_X$ nor $k_\mathrm{Cr}$ can replace these derivations: both would restore the product of drive and fraction but leave a switch operating in a tenth of its range, without its contrast.
The pairing $k_{A,\mathrm{on}}=k_{E,\mathrm{on}}k_{A,\mathrm{off}}/k_{E,\mathrm{off}}$ removes a saturation that the earlier parameterisation had (RasGEFA 83\,\% active at rest and 99\,\% at 60\,nM), where a cAMP step drove RasC-GTP down instead of up; RasC is activated rapidly and transiently by cAMP~\cite{kae2004}.
The pool $[\mathrm{GEF_A}]_\mathrm{tot}$ was chosen so that the number of active GEFs at the driven state is preserved.
The relay is nearly dose independent in the model: 2\,nM cAMP fires the excitable element as fully as 200\,nM, because the LEGI/PIP3 stage reports contrast rather than level; this is why the gate must be steep.
